\begin{pastexam}[label={px:ust-111-ee-ds-05}]{臺聯大電機類 111 學年度 第 5 題}
（問答）

(10\%) Is Heap sort stable? Justify your answer by giving examples.
\end{pastexam}

\begin{pastsoln}
答：不穩定。

穩定的意思是：鍵值相同的元素，排序後的先後順序和排序前一樣（資結書第~14.1~節）。 要說明一個
演算法不穩定，舉一個會把相同鍵值的順序弄反的輸入就夠了。

用資結書虛擬碼~14.13 的堆積排序（建最大堆積，再反覆把根換到最後、縮小堆積、往下調整），排
$\langle 2_a, 2_b, 1 \rangle$，下標只是用來分辨兩個 $2$，原本 $2_a$ 在 $2_b$ 前面。
\begin{enumerate}
\item 建最大堆積：根 $2_a$ 的兩個子節點是 $2_b$、$1$，都不比它大，已經是最大堆積，不動：
$\langle 2_a, 2_b, 1 \rangle$。
\item 根 $2_a$ 和最後一格交換：$\langle 1, 2_b, \mathbf{2_a} \rangle$，堆積剩前兩格。 $1$ 往下調整，
和子節點 $2_b$ 交換：$\langle 2_b, 1, 2_a \rangle$。
\item 根 $2_b$ 和第 2 格交換：$\langle 1, \mathbf{2_b}, 2_a \rangle$，堆積剩一格，結束。
\end{enumerate}
結果是 $\langle 1, 2_b, 2_a \rangle$，兩個 $2$ 的順序反了，所以堆積排序不穩定。

原因：每一輪把根（目前最大的）換到陣列後面，而根是哪一個相同鍵值的元素，取決於堆積裡的位置，
和原本的先後無關；換到後面時還會跳過中間的元素。 上例 $2_a$ 原本在前面，卻先被當成最大值放到
最後一格。

\textbf{另一個例子}：$\langle 3, 1_a, 1_b \rangle$。 建堆積不動；$3$ 換到最後得 $\langle 1_b, 1_a, 3 \rangle$，
堆積 $\langle 1_b, 1_a \rangle$ 不用調整；根 $1_b$ 換到第 2 格得 $\langle 1_a, 1_b, 3 \rangle$。 這次順序
沒變。 可見不穩定指的是「有可能」打亂，不是每次都打亂；只要有一個反例就不穩定。
\end{pastsoln}
